\documentclass{article}
\usepackage{amsmath,amssymb,amsthm}
\begin{document}

\textbf{Subproblem 6 --- Classification of the two-direction case.}

\medskip
\textbf{Statement.}
Let $S\subset\mathbb{R}^2$ be a finite set of $n\ge 3$ points in convex position such that
$P=\operatorname{conv}(S)$ has exactly two distinct edge-direction classes
(i.e.\ every edge direction appears paired with a parallel opposite edge, and there are
exactly two such pairs).

Then $n=4$ and $S$ is the vertex set of a parallelogram.

\medskip
\textbf{Hint.}
In a convex polygon, no two \emph{adjacent} edges are parallel (otherwise two vertices would
coincide or the polygon would be degenerate).
If there are exactly two edge-direction classes, then the edges alternate between the two
directions around the polygon.  An alternating sequence of two directions has even length,
so $n$ is even.
The shortest such polygon has $n=4$.  For $n\ge 6$, the two direction classes each appear
at least twice, and one checks that the only convex polygon (in convex position) with this
property and $n=4$ has opposite sides parallel and equal --- i.e.\ is a parallelogram.
(For $n\ge 6$: two edges of the same direction class that are not adjacent produce a
re-entrant angle or a non-convex arrangement, since the edge vectors must sum to zero around
the polygon and with only two direction classes the sum forces all edges in each class to be
equal and parallel, giving a centrally symmetric polygon; but then $S$ being in convex position
with $n\ge 6$ points from only two edge directions forces repeated vertices, contradiction.)

Conclude: $n=4$ and $P$ is a parallelogram, so $S=\{A,B,C,D\}$ with
$A+C=B+D$ (diagonals bisect each other).

\end{document}
